\subsection*{全体概要：この章で学ぶこと}
\addcontentsline{toc}{subsection}{全体概要：この章で学ぶこと}
この章は、ソフトウェア品質を、製品そのもの、製品を作るプロセス、途中で作られる作業成果物、そして品質を確保する活動の観点から扱う。品質は、要求に合っているかだけでなく、その要求が利害関係者の本当のニーズを表しているかにも依存する。よって、品質を考えるときは、要求、測定、レビュー、テスト、構成管理、標準、倫理、リスクを切り離して扱うことはできない。

到達目標\par
この章を読み終えると、次のことができるようになる。
\begin{itemize}
  \item ソフトウェア品質、製品品質、プロセス品質、品質要求の違いを説明できる。
  \item 誤り、欠陥、故障の違いを説明できる。
  \item 品質コストを、予防、評価、不適合、手戻りの観点から説明できる。
  \item 品質マネジメント、品質保証、品質制御、検証と妥当性確認の役割を区別できる。
  \item 静的解析、動的解析、形式的解析、レビュー、監査、テストの使い分けを説明できる。
  \item 測定、欠陥分類、根本原因分析、是正・予防処置を品質改善に結びつけて説明できる。
\end{itemize}

\par\smallskip
\begin{center}
\centering
\IfFileExists{assets/generated_figures/ch12/fig-ch12-01.png}{\includegraphics[width=0.96\linewidth]{assets/generated_figures/ch12/fig-ch12-01.png}}{\fbox{\parbox[c][48mm][c]{0.9\linewidth}{\centering 画像生成待ち\\12 章の全体地図}}}
\par\smallskip
\refstepcounter{figure}\label{fig:ch12-01}図 \thefigure: 12 章の全体地図
\end{center}
図 \thefigure は、12 章の全体地図を視覚的に整理したものである。
\par\smallskip
